% Shared preamble for all spec documents.
% Include from each spec: % Shared preamble for all spec documents.
% Include from each spec: % Shared preamble for all spec documents.
% Include from each spec: \input{../preamble.tex} (adjust depth).

\documentclass[11pt,a4paper]{article}

\usepackage[utf8]{inputenc}
\usepackage[T1]{fontenc}
\usepackage{lmodern}
\usepackage[a4paper,margin=2.3cm]{geometry}
\usepackage{parskip}
\usepackage{microtype}
\usepackage[dvipsnames]{xcolor}
\usepackage{enumitem}
\usepackage{listings}
\usepackage{tcolorbox}
\usepackage{titlesec}
\usepackage{fancyhdr}
\usepackage{array}
\usepackage{longtable}
\usepackage{booktabs}
\usepackage{amsmath}
\usepackage{amssymb}
\usepackage{hyperref}
\usepackage{cleveref}

\tcbuselibrary{skins,breakable}

\definecolor{specBlue}{RGB}{30,64,132}
\definecolor{specGray}{RGB}{90,90,90}
\definecolor{specAccent}{RGB}{198,40,40}
\definecolor{specGreen}{RGB}{30,110,60}

\hypersetup{
  colorlinks=true,
  linkcolor=specBlue,
  urlcolor=specBlue,
  citecolor=specBlue,
  pdfborder={0 0 0}
}

\titleformat{\section}{\Large\bfseries\color{specBlue}}{\thesection}{0.75em}{}
\titleformat{\subsection}{\large\bfseries\color{specBlue!80}}{\thesubsection}{0.6em}{}

\makeatletter
\newcommand{\specID}[1]{\def\@specID{#1}}
\newcommand{\specTitle}[1]{\def\@specTitle{#1}}
\newcommand{\specStatus}[1]{\def\@specStatus{#1}}
\newcommand{\specVersion}[1]{\def\@specVersion{#1}}
\newcommand{\specOwner}[1]{\def\@specOwner{#1}}
\newcommand{\specTraces}[1]{\def\@specTraces{#1}}

\specID{SPEC-XXX}
\specTitle{Untitled Spec}
\specStatus{DRAFT}
\specVersion{0.1}
\specOwner{Jeremie Nunez}
\specTraces{}

\newcommand{\makespecheader}{%
  \begin{center}
    {\LARGE\bfseries\color{specBlue}\@specTitle}\\[0.4em]
    {\small\color{specGray}\texttt{\@specID} \quad v\@specVersion \quad \textsc{\@specStatus} \quad \@specOwner}
  \end{center}
  \vspace{0.4em}
  \hrule
  \vspace{0.8em}
}

\pagestyle{fancy}
\fancyhf{}
\fancyhead[L]{\small\color{specGray}\@specID}
\fancyhead[R]{\small\color{specGray}\@specTitle}
\fancyfoot[C]{\small\color{specGray}\thepage}
\renewcommand{\headrulewidth}{0pt}
\makeatother

\newtcolorbox{invariant}{
  colback=specBlue!5, colframe=specBlue!75,
  title=Invariant, fonttitle=\bfseries, breakable,
  boxrule=0.5pt, arc=2pt
}

\newtcolorbox{scenario}[1]{
  colback=white, colframe=specBlue!50,
  title={Scenario: #1}, fonttitle=\bfseries, breakable,
  boxrule=0.5pt, arc=2pt
}

\newtcolorbox{ac}[1]{
  colback=specAccent!5, colframe=specAccent!60,
  title={Acceptance Criterion #1}, fonttitle=\bfseries, breakable,
  boxrule=0.5pt, arc=2pt
}

\newtcolorbox{nongoal}{
  colback=specGray!8, colframe=specGray!60,
  title=Non-Goal, fonttitle=\bfseries, breakable,
  boxrule=0.5pt, arc=2pt
}

\lstdefinestyle{codestyle}{
  basicstyle=\ttfamily\small,
  breaklines=true,
  showstringspaces=false,
  frame=single,
  framesep=4pt,
  rulecolor=\color{specBlue!30},
  backgroundcolor=\color{specBlue!2},
  xleftmargin=0.5em,
  xrightmargin=0.5em,
  keywordstyle=\color{specBlue}\bfseries,
  commentstyle=\color{specGray}\itshape,
  stringstyle=\color{specGreen}
}
\lstset{
  inputencoding=utf8,
  extendedchars=true,
  literate=
    {—}{{--}}1
    {–}{{-}}1
    {…}{{...}}1
    {’}{{'}}1
    {‘}{{'}}1
    {“}{{``}}1
    {”}{{''}}1
    {é}{{\'e}}1
    {è}{{\`e}}1
    {ê}{{\^e}}1
    {à}{{\`a}}1
    {â}{{\^a}}1
    {ç}{{\c{c}}}1
    {ô}{{\^o}}1
    {→}{{$\to$}}1
    {←}{{$\leftarrow$}}1
    {×}{{$\times$}}1
    {≥}{{$\geq$}}1
    {≤}{{$\leq$}}1
    {≠}{{$\neq$}}1
    {∈}{{$\in$}}1
    {⌈}{{$\lceil$}}1
    {⌉}{{$\rceil$}}1
    {∞}{{$\infty$}}1
    {α}{{$\alpha$}}1
    {β}{{$\beta$}}1
}
\lstdefinelanguage{TypeScript}{
  keywords={const,let,var,function,return,if,else,for,while,class,interface,type,extends,implements,import,export,from,as,async,await,new,this,true,false,null,undefined,void,any,unknown,never,string,number,boolean,readonly,private,public,protected,enum},
  sensitive=true,
  comment=[l]{//},
  morecomment=[s]{/*}{*/},
  morestring=[b]",
  morestring=[b]',
  morestring=[b]`
}
\lstset{style=codestyle}

\newcommand{\Given}{\textbf{\color{specBlue}Given} }
\newcommand{\When}{\textbf{\color{specBlue}When} }
\newcommand{\Then}{\textbf{\color{specBlue}Then} }
\providecommand{\And}{}
\renewcommand{\And}{\textbf{\color{specBlue}And} }
 (adjust depth).

\documentclass[11pt,a4paper]{article}

\usepackage[utf8]{inputenc}
\usepackage[T1]{fontenc}
\usepackage{lmodern}
\usepackage[a4paper,margin=2.3cm]{geometry}
\usepackage{parskip}
\usepackage{microtype}
\usepackage[dvipsnames]{xcolor}
\usepackage{enumitem}
\usepackage{listings}
\usepackage{tcolorbox}
\usepackage{titlesec}
\usepackage{fancyhdr}
\usepackage{array}
\usepackage{longtable}
\usepackage{booktabs}
\usepackage{amsmath}
\usepackage{amssymb}
\usepackage{hyperref}
\usepackage{cleveref}

\tcbuselibrary{skins,breakable}

\definecolor{specBlue}{RGB}{30,64,132}
\definecolor{specGray}{RGB}{90,90,90}
\definecolor{specAccent}{RGB}{198,40,40}
\definecolor{specGreen}{RGB}{30,110,60}

\hypersetup{
  colorlinks=true,
  linkcolor=specBlue,
  urlcolor=specBlue,
  citecolor=specBlue,
  pdfborder={0 0 0}
}

\titleformat{\section}{\Large\bfseries\color{specBlue}}{\thesection}{0.75em}{}
\titleformat{\subsection}{\large\bfseries\color{specBlue!80}}{\thesubsection}{0.6em}{}

\makeatletter
\newcommand{\specID}[1]{\def\@specID{#1}}
\newcommand{\specTitle}[1]{\def\@specTitle{#1}}
\newcommand{\specStatus}[1]{\def\@specStatus{#1}}
\newcommand{\specVersion}[1]{\def\@specVersion{#1}}
\newcommand{\specOwner}[1]{\def\@specOwner{#1}}
\newcommand{\specTraces}[1]{\def\@specTraces{#1}}

\specID{SPEC-XXX}
\specTitle{Untitled Spec}
\specStatus{DRAFT}
\specVersion{0.1}
\specOwner{Jeremie Nunez}
\specTraces{}

\newcommand{\makespecheader}{%
  \begin{center}
    {\LARGE\bfseries\color{specBlue}\@specTitle}\\[0.4em]
    {\small\color{specGray}\texttt{\@specID} \quad v\@specVersion \quad \textsc{\@specStatus} \quad \@specOwner}
  \end{center}
  \vspace{0.4em}
  \hrule
  \vspace{0.8em}
}

\pagestyle{fancy}
\fancyhf{}
\fancyhead[L]{\small\color{specGray}\@specID}
\fancyhead[R]{\small\color{specGray}\@specTitle}
\fancyfoot[C]{\small\color{specGray}\thepage}
\renewcommand{\headrulewidth}{0pt}
\makeatother

\newtcolorbox{invariant}{
  colback=specBlue!5, colframe=specBlue!75,
  title=Invariant, fonttitle=\bfseries, breakable,
  boxrule=0.5pt, arc=2pt
}

\newtcolorbox{scenario}[1]{
  colback=white, colframe=specBlue!50,
  title={Scenario: #1}, fonttitle=\bfseries, breakable,
  boxrule=0.5pt, arc=2pt
}

\newtcolorbox{ac}[1]{
  colback=specAccent!5, colframe=specAccent!60,
  title={Acceptance Criterion #1}, fonttitle=\bfseries, breakable,
  boxrule=0.5pt, arc=2pt
}

\newtcolorbox{nongoal}{
  colback=specGray!8, colframe=specGray!60,
  title=Non-Goal, fonttitle=\bfseries, breakable,
  boxrule=0.5pt, arc=2pt
}

\lstdefinestyle{codestyle}{
  basicstyle=\ttfamily\small,
  breaklines=true,
  showstringspaces=false,
  frame=single,
  framesep=4pt,
  rulecolor=\color{specBlue!30},
  backgroundcolor=\color{specBlue!2},
  xleftmargin=0.5em,
  xrightmargin=0.5em,
  keywordstyle=\color{specBlue}\bfseries,
  commentstyle=\color{specGray}\itshape,
  stringstyle=\color{specGreen}
}
\lstset{
  inputencoding=utf8,
  extendedchars=true,
  literate=
    {—}{{--}}1
    {–}{{-}}1
    {…}{{...}}1
    {’}{{'}}1
    {‘}{{'}}1
    {“}{{``}}1
    {”}{{''}}1
    {é}{{\'e}}1
    {è}{{\`e}}1
    {ê}{{\^e}}1
    {à}{{\`a}}1
    {â}{{\^a}}1
    {ç}{{\c{c}}}1
    {ô}{{\^o}}1
    {→}{{$\to$}}1
    {←}{{$\leftarrow$}}1
    {×}{{$\times$}}1
    {≥}{{$\geq$}}1
    {≤}{{$\leq$}}1
    {≠}{{$\neq$}}1
    {∈}{{$\in$}}1
    {⌈}{{$\lceil$}}1
    {⌉}{{$\rceil$}}1
    {∞}{{$\infty$}}1
    {α}{{$\alpha$}}1
    {β}{{$\beta$}}1
}
\lstdefinelanguage{TypeScript}{
  keywords={const,let,var,function,return,if,else,for,while,class,interface,type,extends,implements,import,export,from,as,async,await,new,this,true,false,null,undefined,void,any,unknown,never,string,number,boolean,readonly,private,public,protected,enum},
  sensitive=true,
  comment=[l]{//},
  morecomment=[s]{/*}{*/},
  morestring=[b]",
  morestring=[b]',
  morestring=[b]`
}
\lstset{style=codestyle}

\newcommand{\Given}{\textbf{\color{specBlue}Given} }
\newcommand{\When}{\textbf{\color{specBlue}When} }
\newcommand{\Then}{\textbf{\color{specBlue}Then} }
\providecommand{\And}{}
\renewcommand{\And}{\textbf{\color{specBlue}And} }
 (adjust depth).

\documentclass[11pt,a4paper]{article}

\usepackage[utf8]{inputenc}
\usepackage[T1]{fontenc}
\usepackage{lmodern}
\usepackage[a4paper,margin=2.3cm]{geometry}
\usepackage{parskip}
\usepackage{microtype}
\usepackage[dvipsnames]{xcolor}
\usepackage{enumitem}
\usepackage{listings}
\usepackage{tcolorbox}
\usepackage{titlesec}
\usepackage{fancyhdr}
\usepackage{array}
\usepackage{longtable}
\usepackage{booktabs}
\usepackage{amsmath}
\usepackage{amssymb}
\usepackage{hyperref}
\usepackage{cleveref}

\tcbuselibrary{skins,breakable}

\definecolor{specBlue}{RGB}{30,64,132}
\definecolor{specGray}{RGB}{90,90,90}
\definecolor{specAccent}{RGB}{198,40,40}
\definecolor{specGreen}{RGB}{30,110,60}

\hypersetup{
  colorlinks=true,
  linkcolor=specBlue,
  urlcolor=specBlue,
  citecolor=specBlue,
  pdfborder={0 0 0}
}

\titleformat{\section}{\Large\bfseries\color{specBlue}}{\thesection}{0.75em}{}
\titleformat{\subsection}{\large\bfseries\color{specBlue!80}}{\thesubsection}{0.6em}{}

\makeatletter
\newcommand{\specID}[1]{\def\@specID{#1}}
\newcommand{\specTitle}[1]{\def\@specTitle{#1}}
\newcommand{\specStatus}[1]{\def\@specStatus{#1}}
\newcommand{\specVersion}[1]{\def\@specVersion{#1}}
\newcommand{\specOwner}[1]{\def\@specOwner{#1}}
\newcommand{\specTraces}[1]{\def\@specTraces{#1}}

\specID{SPEC-XXX}
\specTitle{Untitled Spec}
\specStatus{DRAFT}
\specVersion{0.1}
\specOwner{Jeremie Nunez}
\specTraces{}

\newcommand{\makespecheader}{%
  \begin{center}
    {\LARGE\bfseries\color{specBlue}\@specTitle}\\[0.4em]
    {\small\color{specGray}\texttt{\@specID} \quad v\@specVersion \quad \textsc{\@specStatus} \quad \@specOwner}
  \end{center}
  \vspace{0.4em}
  \hrule
  \vspace{0.8em}
}

\pagestyle{fancy}
\fancyhf{}
\fancyhead[L]{\small\color{specGray}\@specID}
\fancyhead[R]{\small\color{specGray}\@specTitle}
\fancyfoot[C]{\small\color{specGray}\thepage}
\renewcommand{\headrulewidth}{0pt}
\makeatother

\newtcolorbox{invariant}{
  colback=specBlue!5, colframe=specBlue!75,
  title=Invariant, fonttitle=\bfseries, breakable,
  boxrule=0.5pt, arc=2pt
}

\newtcolorbox{scenario}[1]{
  colback=white, colframe=specBlue!50,
  title={Scenario: #1}, fonttitle=\bfseries, breakable,
  boxrule=0.5pt, arc=2pt
}

\newtcolorbox{ac}[1]{
  colback=specAccent!5, colframe=specAccent!60,
  title={Acceptance Criterion #1}, fonttitle=\bfseries, breakable,
  boxrule=0.5pt, arc=2pt
}

\newtcolorbox{nongoal}{
  colback=specGray!8, colframe=specGray!60,
  title=Non-Goal, fonttitle=\bfseries, breakable,
  boxrule=0.5pt, arc=2pt
}

\lstdefinestyle{codestyle}{
  basicstyle=\ttfamily\small,
  breaklines=true,
  showstringspaces=false,
  frame=single,
  framesep=4pt,
  rulecolor=\color{specBlue!30},
  backgroundcolor=\color{specBlue!2},
  xleftmargin=0.5em,
  xrightmargin=0.5em,
  keywordstyle=\color{specBlue}\bfseries,
  commentstyle=\color{specGray}\itshape,
  stringstyle=\color{specGreen}
}
\lstset{
  inputencoding=utf8,
  extendedchars=true,
  literate=
    {—}{{--}}1
    {–}{{-}}1
    {…}{{...}}1
    {’}{{'}}1
    {‘}{{'}}1
    {“}{{``}}1
    {”}{{''}}1
    {é}{{\'e}}1
    {è}{{\`e}}1
    {ê}{{\^e}}1
    {à}{{\`a}}1
    {â}{{\^a}}1
    {ç}{{\c{c}}}1
    {ô}{{\^o}}1
    {→}{{$\to$}}1
    {←}{{$\leftarrow$}}1
    {×}{{$\times$}}1
    {≥}{{$\geq$}}1
    {≤}{{$\leq$}}1
    {≠}{{$\neq$}}1
    {∈}{{$\in$}}1
    {⌈}{{$\lceil$}}1
    {⌉}{{$\rceil$}}1
    {∞}{{$\infty$}}1
    {α}{{$\alpha$}}1
    {β}{{$\beta$}}1
}
\lstdefinelanguage{TypeScript}{
  keywords={const,let,var,function,return,if,else,for,while,class,interface,type,extends,implements,import,export,from,as,async,await,new,this,true,false,null,undefined,void,any,unknown,never,string,number,boolean,readonly,private,public,protected,enum},
  sensitive=true,
  comment=[l]{//},
  morecomment=[s]{/*}{*/},
  morestring=[b]",
  morestring=[b]',
  morestring=[b]`
}
\lstset{style=codestyle}

\newcommand{\Given}{\textbf{\color{specBlue}Given} }
\newcommand{\When}{\textbf{\color{specBlue}When} }
\newcommand{\Then}{\textbf{\color{specBlue}Then} }
\providecommand{\And}{}
\renewcommand{\And}{\textbf{\color{specBlue}And} }


\specID{SPEC-SW-010}
\specTitle{Feedback Loop --- reviewer signals tune next swarm run per category}
\specStatus{DRAFT}
\specVersion{0.1}
\specOwner{Jeremie Nunez}

\begin{document}
\makespecheader

\section{Context}
The sweep subsystem runs a fan-out of cheap nano workers over the domain catalogue, each producing
\emph{suggestions} that a human reviewer either accepts or rejects. Without a feedback channel, the
swarm repeats the same low-signal suggestions cycle after cycle: the reviewer ends up rejecting the
same kind of false positive every week. This spec defines the closed loop that turns reviewer
decisions into per-category prompt conditioning for the next swarm run, so rejection rates decay
over time without any manual prompt editing.

\section{Goals}
\begin{itemize}[leftmargin=1.2em]
  \item Persist every accept/reject decision with its category, target entity, and optional reviewer note.
  \item Inject the most recent decisions back into the nano prompt for the \emph{same} category on the next run, within one sweep cycle.
  \item Make the effect observable: per-category rejection rate must be queryable as a time series and must trend down when the signal is informative.
  \item Bound the feedback block size so prompt tokens stay within the configured budget regardless of feedback volume.
  \item Fail open: if the feedback store is unreachable, the swarm runs without conditioning rather than blocking.
\end{itemize}

\section{Non-Goals}
\begin{nongoal}
This spec does not cover: (a) how suggestions are generated (SPEC-SW-001 Nano Sweep), (b) the
reviewer UI surface, (c) cross-category feedback (a reject in \texttt{price\_anomaly} does not
condition \texttt{missing\_data}), (d) model fine-tuning or embedding-based recall of past
feedback --- only literal injection of recent decisions is in scope, (e) long-term feedback
archiving beyond the rolling window.
\end{nongoal}

\section{Invariants}
\begin{invariant}
A reject decision recorded before sweep cycle $N$ starts \emph{must} be reflected in the prompt
used for its category during cycle $N$. There is no grace cycle: one cycle of latency is the
upper bound, not a target.
\end{invariant}

\begin{invariant}
Feedback is category-scoped: the prompt for category $C$ in cycle $N$ only embeds decisions
whose stored category equals $C$. No cross-leakage between categories.
\end{invariant}

\begin{invariant}
The injected feedback block is bounded: at most \texttt{MAX\_FEEDBACK\_LINES} entries per category
per batch, ordered most-recent-first, truncated by reviewer note length to fit \texttt{MAX\_FEEDBACK\_CHARS}.
\end{invariant}

\begin{invariant}
Feedback store unavailability degrades gracefully: the sweep still runs, with an empty feedback
block, and emits a warning log --- it never throws into the caller.
\end{invariant}

\begin{invariant}
Per-category rejection rate is a monotonically defined observable: for any window $[t_0, t_1]$ and
category $C$, \texttt{rejectionRate}$(C, t_0, t_1) = \texttt{rejected}(C) / (\texttt{rejected}(C) + \texttt{accepted}(C))$,
computed from persisted decisions only.
\end{invariant}

\section{Interface}
Public surface exposed by the feedback subsystem. Consumers (nano sweep service, admin
dashboards) depend on these signatures; internal storage layout is free to change.

\begin{lstlisting}[language=TypeScript]
// Persisted decision record
export interface ReviewerDecision {
  id: string;
  category: SweepCategory;
  entityName: string | null;     // e.g. appellation / dataset / segment
  wasAccepted: boolean;
  reviewerNote: string | null;
  decidedAt: Date;
}

// Loaded into the prompt builder
export interface PastFeedback {
  category: SweepCategory;
  entityName: string | null;
  wasAccepted: boolean;
  reviewerNote: string | null;
}

export interface FeedbackRepository {
  recordDecision(input: Omit<ReviewerDecision, "id" | "decidedAt">): Promise<ReviewerDecision>;
  loadPastFeedback(opts?: {
    category?: SweepCategory;
    limit?: number;          // defaults to MAX_FEEDBACK_LINES
  }): Promise<PastFeedback[]>;
  rejectionRate(
    category: SweepCategory,
    window: { from: Date; to: Date },
  ): Promise<{ rejected: number; accepted: number; rate: number }>;
}

// Prompt builder contract — pure function, easy to test.
export function buildFeedbackBlock(
  category: SweepCategory,
  feedback: PastFeedback[],
): string;
\end{lstlisting}

\section{Scenarios}

\begin{scenario}{Single reject is reflected in the next cycle}
\Given the reviewer rejects a suggestion in category \texttt{price\_anomaly} for entity \texttt{E} with note \texttt{"price is within the historical band"} at $t_0$ \\
\And no sweep is currently in flight \\
\When a new sweep cycle starts at $t_1 > t_0$ \\
\Then the nano prompt generated for any batch containing entity \texttt{E} in category \texttt{price\_anomaly} contains a feedback line referencing \texttt{E}, marked \texttt{REJECTED}, with the reviewer note verbatim
\end{scenario}

\begin{scenario}{Category isolation}
\Given a reject recorded for category \texttt{missing\_data} on entity \texttt{E} \\
\When the next sweep runs and produces a batch targeting entity \texttt{E} \\
\Then the prompt for category \texttt{price\_anomaly} on that batch contains no feedback line for \texttt{E} \\
\And the prompt for category \texttt{missing\_data} on that batch does contain the feedback line
\end{scenario}

\begin{scenario}{Feedback block is bounded}
\Given 500 rejects exist for category \texttt{enrichment} \\
\When the next sweep cycle builds a prompt for that category \\
\Then the injected feedback block contains at most \texttt{MAX\_FEEDBACK\_LINES} lines \\
\And the entries are ordered most-recent-first \\
\And the serialized block size does not exceed \texttt{MAX\_FEEDBACK\_CHARS} characters
\end{scenario}

\begin{scenario}{Feedback store unavailable}
\Given the feedback repository raises on \texttt{loadPastFeedback} \\
\When a sweep cycle starts \\
\Then the sweep completes with an empty feedback block \\
\And a warning is logged with the underlying error \\
\And no exception escapes the service boundary
\end{scenario}

\begin{scenario}{Rejection rate decay is observable}
\Given 4 consecutive weekly sweep cycles, each preceded by at least one reject in category \texttt{price\_anomaly} with informative reviewer notes \\
\When the rejection rate is computed for each weekly window \\
\Then the rate for week $N+1$ is less than or equal to the rate for week $N$, within a documented tolerance, for categories where reviewer notes were informative \\
\And the raw \texttt{rejected} and \texttt{accepted} counts are exposed per window
\end{scenario}

\section{Acceptance Criteria}

\begin{ac}{AC-1}
\texttt{recordDecision} persists all fields; a subsequent \texttt{loadPastFeedback} with a matching category returns the record.
\end{ac}

\begin{ac}{AC-2}
For a reject recorded at $t_0$, the first sweep cycle started after $t_0$ produces, for the matching category and entity, a prompt whose feedback block contains a line tagged \texttt{REJECTED} with the reviewer note.
\end{ac}

\begin{ac}{AC-3}
Given 1000 stored decisions across all categories, \texttt{loadPastFeedback(\{category: C\})} returns at most \texttt{MAX\_FEEDBACK\_LINES} entries, all with category equal to $C$.
\end{ac}

\begin{ac}{AC-4}
\texttt{buildFeedbackBlock(C, feedback)} emits no line belonging to a category other than $C$, given mixed-category input.
\end{ac}

\begin{ac}{AC-5}
When the feedback repository throws, the sweep service completes successfully with a non-empty result and an empty feedback block; a warning log line is emitted.
\end{ac}

\begin{ac}{AC-6}
\texttt{rejectionRate(C, window)} returns values consistent with $\texttt{rate} = \texttt{rejected} / (\texttt{rejected} + \texttt{accepted})$ and $0$ when the window contains no decisions.
\end{ac}

\begin{ac}{AC-7}
In a fixture simulating 4 cycles with informative rejects for category $C$, the computed rejection rate for cycle 4 is strictly less than that of cycle 1.
\end{ac}

\section{Traceability}

\begin{center}
\begin{tabular}{@{}lll@{}}
\toprule
AC & Test file & Test name \\
\midrule
AC-1 & \texttt{packages/sweep/tests/unit/feedback-repository.spec.ts} & \texttt{records and reloads a decision} \\
AC-2 & \texttt{packages/sweep/tests/integration/feedback-loop.spec.ts} & \texttt{reject surfaces in next cycle prompt} \\
AC-3 & \texttt{packages/sweep/tests/unit/feedback-repository.spec.ts} & \texttt{loadPastFeedback caps at MAX\_FEEDBACK\_LINES} \\
AC-4 & \texttt{packages/sweep/tests/unit/feedback-prompt.spec.ts} & \texttt{buildFeedbackBlock is category-scoped} \\
AC-5 & \texttt{packages/sweep/tests/integration/feedback-loop.spec.ts} & \texttt{repository failure degrades gracefully} \\
AC-6 & \texttt{packages/sweep/tests/unit/feedback-repository.spec.ts} & \texttt{rejectionRate math} \\
AC-7 & \texttt{packages/sweep/tests/integration/feedback-decay.spec.ts} & \texttt{rejection rate decays across cycles} \\
\bottomrule
\end{tabular}
\end{center}

\section{Open Questions}
\begin{itemize}[leftmargin=1.2em]
  \item Should accepted decisions also be injected, or only rejects? Current draft injects both, tagged accordingly.
  \item What is the right default for \texttt{MAX\_FEEDBACK\_LINES} (current proposal: 5)?
  \item Do we need an explicit \texttt{expiresAt} on decisions, or is the rolling most-recent window sufficient?
  \item Should reviewer notes be truncated per-entry, or the block truncated as a whole once the char budget is hit?
\end{itemize}

\end{document}
